% New standalone increment. Prior increments and GitHub remain read-only.
\documentclass[11pt]{article}
\usepackage[T1]{fontenc}\usepackage[utf8]{inputenc}\usepackage{lmodern}
\usepackage[margin=.87in]{geometry}
\usepackage{amsmath,amssymb,amsthm,mathtools,microtype,array,longtable,booktabs,xurl}
\usepackage[hidelinks]{hyperref}
\setlength{\emergencystretch}{3em}\allowdisplaybreaks[2]
\numberwithin{equation}{section}
\newtheorem{theorem}{Theorem}[section]
\newtheorem{lemma}[theorem]{Lemma}
\newtheorem{proposition}[theorem]{Proposition}
\newtheorem{corollary}[theorem]{Corollary}
\newcommand{\E}{\mathbb E}\newcommand{\dd}{\,d}
\newcommand{\KK}{\mathcal K}\newcommand{\BB}{\mathcal B}
\newcommand{\Gp}{G_+}\newcommand{\dist}{\mathfrak d}
\newcommand{\one}{\mathbf1}\newcommand{\Kfun}{\mathsf K}
\newcommand{\Ffun}{\mathsf F}
\title{Source-rate transfer and initialized learned reversal\\
with persistence in the deep regime}
\author{Standalone analytical increment for independent review}
\date{Version 1 --- 2 November 2026}
\begin{document}\maketitle
\begin{abstract}
FWC's finite-window continuation is combined with WIT's learned
limiting witnesses. A pairwise core comparison supplies a strict
probe gap for every core label. The exact original source rates
then differ from their leading-form counterparts by a vanishing
normalized error, including all background and radial-reaction terms.
Both concept responses and the fixed cluster-loss reduction are
proved directly, without assuming compact weak support or deleting
outside labels. This gives initialized Theorem A on WIT's nonempty
parameter rectangle for \(S_0=q^{10}\), with source signs,
ordered total-rate windows and a quantitative rebound.
FWC's every-fixed-window error transfer also discharges LP's
initialized persistence premise. No quantitative growing horizon,
canonical-regime assertion or original-flow lag decomposition is made.
\end{abstract}

\section{Scope, clocks, exact conventions and hypothesis inventory}
\label{an07rtav1:inventory}
Keep the original balanced untied population initialized by the
uniform label law with \(m(0,b)=S_0=q^{10}\).
In normalized coordinates
\[
 P_1=N(e_1,q^2I_2),\qquad P_2=N(\rho e_2,q^2I_2),\qquad
 P_{\rm mix}=\tfrac12(P_1+P_2),\qquad \lambda=\rho^2.
\]
Write \(U=\sqrt m\,a_\theta\), \(W=\sqrt m\,a_\psi\),
\(a_\theta=(\cos\theta,\sin\theta)\), and
\(d\mu=m\,d\lambda_0\). All cohorts are the unchanged fixed
initial-label sets \(\KK,\Gp,\Sigma,\BB\), with
\(S=\mathcal C_q\subset\Gp\), \(P_q=\Gp\setminus S\).
Use \(M=\mu_\KK\) and \(N=\mu_S\) for physical masses.
The output and probe error are
\[
 f(X)=\int m a_\psi(a_\theta\cdot X)_+\dd\lambda_0,
 \quad \xi_q(\zeta)=(\sin(q\zeta),-\cos(q\zeta)),
 \quad E_{\xi_q}=\tfrac12|\xi_q-f(\xi_q)|^2.
\]
Source matrices include the mixture half:
\begin{equation}
 R^{(p)}_\theta=\tfrac12\E_{P_p}[(X-f(X))X^T
                    \one_{\{U\cdot X>0\}}],\qquad
 U'_p=(R^{(p)})^TW,\quad W'_p=R^{(p)}U.
 \label{an07rtav1:sourceconvention}
\end{equation}
Set \(h_p=(\log m)'_p=2a_\psi^TR^{(p)}a_\theta\);
this is twice CWI's auxiliary scalar \(g_p\).
The original source rate is
\(D_p^\tau=-(\xi_q-f)\cdot\dot f_p\), with
\(D^\tau=D_1^\tau+D_2^\tau\). At a probe gate use CWI's
selector \(\one_{\{a_\theta\cdot\xi_q>0\}}\).
These are classical source derivatives at gate-regular states and
give the time derivative almost everywhere. No classical derivative
at a background gate-hit atom is presumed below.

Retain the analytically defined WCS/PSE/FWC center
\begin{equation}
 \tau=\mu_1^2t_{\rm phys},\qquad
 \tau_h=t_n+\lambda^{-1}\log\frac{1-n}{n},\qquad
 s=\tau-\tau_h,\qquad n_\rho(s)=(1+e^{-\lambda s})^{-1}.
 \label{an07rtav1:time}
\end{equation}
Here \(n\le1/8\) is the unchanged WCS level and \(t_n\)
is its exact auxiliary-logistic hit. This center matches the
selected branch's half-mass phase; it is not declared to be the
exact half-mass time of the physical \(N\). FWC bounds the
physical mass/logit error on every fixed window. Thus WIT's
\(s\)-clock is exactly the comparison clock in
\eqref{an07rtav1:time}, including its \(\rho\)-dependence.

Keep WIT's \emph{original}, analytically specified rectangle
and witnesses, without recomputing its widths:
\begin{gather}
 \mathcal P_*=K_\rho\times I_\zeta,\quad
 K_\rho=[2/3-\delta_\rho,2/3+\delta_\rho],\quad
 I_\zeta=[5\pi/9-\delta_\zeta,5\pi/9+\delta_\zeta],
 \label{an07rtav1:rectangle}\\
 J=[s_*-r,s_*+r],\quad
 J_-=[s_*-r,s_*-r/2],\quad
 J_+=[s_*+r/2,s_*+r].
 \label{an07rtav1:windows}
\end{gather}
The constants \(\delta_\rho,\delta_\zeta,r,p_*,\gamma_0,\sigma,g_*\)
are positive and defined by WIT's
\texttt{rdef,\allowbreak windows,\allowbreak constants,\allowbreak rectangle}. In particular
\(\mathcal P_*\subset[13/20,7/10]\times[87/50,7/4]\).
WIT \texttt{uniformtheorem} gives on this same rectangle
\begin{align}
 n_\rho&\ge1/4+\sigma/2,& \zeta-X_\rho&\ge g_*/4&&\text{on }J,
 \notag\\
 d_{1,*}&\le-5p_*/8,&d_{2,*}&\ge5p_*/8&&\text{on }J,
 \label{an07rtav1:limitingmargins}\\
 \mathcal E_*'&\le-3\gamma_0/4&&&\text{on }J_-,\notag\\
 \mathcal E_*'&\ge3\gamma_0/4&&&\text{on }J_+.\notag
\end{align}
Since \(n<1/4\), monotonicity of \(n_\rho\) proves
\(s_n(\rho)<s_*-r\) for every \(\rho\in K_\rho\).
Choose one fixed
\(T_J>\max\{s_*+r,\max_{[13/20,7/10]}s_n\}\).
FWC therefore supplies all its conclusions throughout \(J\).
Write \(L=1+\log(1/q)\), \(\chi=1/20\),
\(\gamma_*=289169/924000\) and
\(\epsilon_{\rm bg}=q^8+L^2q^{\gamma_*}\).

\small
\begin{longtable}{>{\raggedright\arraybackslash}p{.21\textwidth}
 >{\raggedright\arraybackslash}p{.43\textwidth}
 >{\raggedright\arraybackslash}p{.27\textwidth}}
\toprule Source & Hypotheses, including domains & Use and supply\\\midrule\endhead
FWC \texttt{main,\allowbreak corefield,\allowbreak distbound,\allowbreak background}
& Original initialized \(\kappa=10\), deep ratio box; inherited
 fixed constants; fixed finite \(T>\max s_n\), bounded nonnegative
 offsets; \(q\) last. Current measures and fixed cohorts.
& Unconditional original-flow conclusion on \([t_n,\tau_h+T_J]\).
 No passage guard is added. Later fixed horizons are used separately
 for persistence.\\
WIT \texttt{uniformtheorem}
& Selected coherent branch, leading strong mass one, exact
 half-mass-centered \(n_\rho\); parameters in its \(\mathcal P_*\);
 its specified common \(J,J_-,J_+\) and positive analytic constants.
& Supplies \eqref{an07rtav1:limitingmargins}. The original physical
 mass is compared by FWC, not set equal to \(n_\rho\).\\
WIT \texttt{populationtransfer,\allowbreak epsdef}
& A solution of the defined leading system at the same \(\rho,s\);
 positive measures, \(M,N\le2\), both supports in \([-R_W,R_W]^2\)
 throughout \(J\), and \(\sup_J\dist\le\varepsilon_*\), with
 \(\varepsilon_*=\min\{1/4,\sigma/4,p_*/(8C_1),\gamma_0/(8C_1)\}\),
 \(C_1=64R_W^2(1+4/g_*)\); selector/a.e. chain rule.
& Not applied to the perturbed original flow or its noncompact weak
 chart. Source rates and learning are transferred directly below.
 Its tolerance can also be met by FWC for small \(q\).\\
WIT \texttt{staticproposition}
& Entire exact state consists of two charted measures with both
 supports in \([-R_W,R_W]^2\), \(R_W\ge2\), \(M,N\le2\),
 \(\rho\in[13/20,7/10]\), \(0<q\le1\). Its final learning claim
 also needs its population mass conclusions, \(S_0\le1/2\),
 and \(q\le q_{\rm learn}\), the minimum of
 \(1,\sigma/(100R_W),1/(100R_W)\) and
 \((13/20)/(272R_W)\).
& Entire-state compactness is not supplied. Section 4 replaces the
 estimate by a direct full-output \(L^2\) proof, including weak tails
 and all omitted labels. No learning input remains open.\\
CWI \texttt{rate}
& Leading-form rate evaluated at positive measures; \(0<\rho\le1\),
 \(M,N\le2\), \(0\le n_*\le1\), \(R\ge2\); strong support and
 comparison coordinates in \([-R,R]^2\); finite weak first moments;
 \(|\zeta|\le R\), \(\zeta-X_*\ge g_0>0\).
& FWC supplies the state and moments, WIT gives \(g_0=g_*/4\).
 This is a static leading-form comparison, not an exact original
 leading-system identification.\\
CWI \texttt{exactrate,\allowbreak ratecharge}; PSE matrix/energy
& Exact balanced output; fixed cohorts; sourcewise bounded
 diagonal and \(O(q)\) cross entries; \(O(qL)\) alignment;
 bounded offsets, selector/a.e. convention.
& FWC continuation and background prove these on \(J\), including
 every strip and angle slice. Its primitive estimates replace the
 old passage-only scope.\\
P \texttt{probe-pop,\allowbreak atomicerror}
& Full population probe expansion assumes two compact charted
 families and no outside; its simple source identities assume
 leading \(M=1\). Atomic formula is for the selected coherent state.
& Only the core expansion is re-used, with its actual reactions;
 background is separately charged. \(M-1\) is retained. Atomic
 formula identifies the limiting observable and clearing.\\
LP \texttt{uniform-clear}
& Compact \(K\subset(0,1)\), selected branch continuous in \(\rho\)
 on finite windows, uniform forward growth bound, finite upper offset.
& Use \(K=K_\rho\), \(\zeta_+=\max I_\zeta\). WIT
 \texttt{branchsmooth} and P's growth inequality supply these,
 as explicitly checked in LPR.\\
LPR \texttt{combinedtheorem}
& Same selected leading branch and WIT rectangle/witnesses;
 imported P and LP continuation, invariants and uniform clearing.
& Supplies one fixed \(S>s_*+r\) with \(\mathcal E_*=0\) for
 all later times and all \(\mathcal P_*\). No original-flow
 conclusion is borrowed from that theorem.\\
LP \texttt{diagonal-transfer}
& Actual initialized positive-noise family; analytical physical
 centering; same compact offsets/parameters and common clearing
 time; full original compact-window convergence for every finite
 \([S,T]\).
& \(S_0=q^{10}\), center \(\tau_h/\mu_1^2\), LPR clearing,
 and FWC \texttt{errortransfer,\allowbreak early} supply all inputs. No rate
 of the growing horizon is inferred.\\
\bottomrule
\end{longtable}\normalsize
The source \(R_W\) in WIT's rectangle/tolerance is frozen.
The possibly larger continuation radius \(R\) used below does not
redefine \(\delta_\rho,\delta_\zeta\) or the common windows.

\section{Every core label is separated from the probe gate}
\label{an07rtav1:coresection}
Choose \(R\ge2\) covering FWC's core chart, all selected comparator
coordinates on \([s_n,T_J]\), and the offsets. Let \(D_J\) be an
upper bound for \(T_J-s_n(\rho)\). All constants in this and the
next sections are fixed before \(q\), uniformly over \(K_\rho\).

\begin{lemma}[Supremum core tracking and strict probe gap]
\label{an07rtav1:coregap}
For some fixed \(C_{\rm sup}\),
\begin{equation}
 \sup_{b\in\KK}|(x_b,y_b)-(X_*,Y_*)|_1
                  \le C_{\rm sup}q^\chi
 \quad\text{on }[t_n,\tau_h+T_J].
 \label{an07rtav1:suptracking}
\end{equation}
In particular, if \(C_{\rm sup}q^\chi\le g_*/8\), then
\(\zeta-x_b\ge g_*/8\) for every core label and every
\((\rho,\zeta,s)\in\mathcal P_*\times J\).
All core labels are strictly probe-active for small \(q\), including
closed endpoints of their label set.
\end{lemma}
\begin{proof}
At fixed time, hold the physical donor measures and moments fixed
in P's leading strong field \(F_s(z;\mu_\KK,\mu_S)\).
Its receiver Lipschitz constant on \([-R,R]^2\) is bounded by a
fixed \(L_F\). Indeed \(M,N\le2\), the current weak first moments
are bounded by FWC, and
\[
 \partial_1\Ffun(t,w)=(w-t)_+\phi(t),\qquad
 |\Kfun'|\le1
\]
bound the strong donor-kernel derivative on this compact chart.
All other receiver derivatives are bounded Gaussian functions times
these bounded moments. No derivative estimate for the original
finite-noise remainder is needed.

FWC \texttt{corefield} gives
\(z_b'=F_s(z_b;\mu_\KK,\mu_S)+e_b\) with
\(\sup_b|e_b|_1\le C_Fq^{1/10}\). Hence any two fixed core labels
satisfy
\[
 D^+|z_b-z_{b'}|_1\le L_F|z_b-z_{b'}|_1+2C_Fq^{1/10}.
\]
The entry core radius is \(O(q^\chi)\) by the WCS/PSE handoff.
Gronwall bounds the entire diameter by
\(e^{L_FD_J}[C_{\rm ent}q^\chi+2C_FD_Jq^{1/10}]=O(q^\chi)\).
For every label, averaging over the current core measure yields
\[
 |z_b-Z_*|_1\le\operatorname{diam}(\KK)+
       M^{-1}\int_\KK|z_{b'}-Z_*|_1\dd\mu(b').
\]
FWC gives \(M\ge1/2\) and the last integral \(O(q^\chi)\).
This proves \eqref{an07rtav1:suptracking} without needing to
identify the chosen reference label with a core endpoint.
Combine with WIT's gap \(\zeta-X_*\ge g_*/4\).
Finally \(a_\theta\cdot\xi_q=\sin(q(\zeta-x_b))>0\), since
the gap is positive and its upper bound is fixed; take
\(2qR<\pi/2\). Continuity extends the estimate to closed labels.
\end{proof}

\section{Exact source rates, background and leading comparison}
\label{an07rtav1:ratesection}
Let \(z=a_\theta\cdot\xi_q\), and let
\(v_{\theta,p},v_{\psi,p}\) be the exact source angular velocities
from \eqref{an07rtav1:sourceconvention}. For any fixed cohort \(A\),
\begin{equation}
 \dot f_{A,p}=\int_A m\big[
 h_p a_\psi z_++v_{\psi,p}a_\psi^\perp z_+
 +v_{\theta,p}a_\psi\one_{\{z>0\}}
                       (a_\theta^\perp\cdot\xi_q)\big]\dd\lambda_0.
 \label{an07rtav1:exactrate}
\end{equation}
There is no extra factor \(1/2\) to apply to this formula.
The primary family is strictly inactive by FWC, so its direct
output and direct source derivatives vanish exactly. Its contribution
to the training residual is retained in all core velocities.

\begin{lemma}[Core self-rate through physical leading-form rates]
\label{an07rtav1:corerate}
Let \(d_p^{\rm lead}\) be the leading-form source rate evaluated
at the actual physical core and selected measures. Uniformly on
\(\mathcal P_*\times J\),
\begin{equation}
 \left|q^{-2}D_p^{\KK,\rm self}-d_p^{\rm lead}\right|
 \le C_{\rm core}(q^{1/10}+q^2),\qquad
 D_p^{\KK,\rm self}=-(\xi_q-f_\KK)\cdot\dot f_{\KK,p}.
 \label{an07rtav1:coreprecision}
\end{equation}
The core velocities in this self-output expression are driven by
the full original field, not by a core-only training problem.
\end{lemma}
\begin{proof}
All core gates are strictly active by Lemma \ref{an07rtav1:coregap}.
Set
\[
 A=\int_\KK(\zeta-x)\dd\mu,\qquad
 B=\int_\KK y(\zeta-x)\dd\mu,
\]
and, with the actual source velocities and full source reaction,
\begin{align*}
 A_p^{\rm act}&=\int_\KK[(\zeta-x)h_p-x'_p]\dd\mu,\\
 B_p^{\rm act}&=\int_\KK[y(\zeta-x)h_p
                  +(\zeta-x)y'_p-yx'_p]\dd\mu.
\end{align*}
Uniform Taylor expansion of the exact active sine and output
coordinates, and of \eqref{an07rtav1:exactrate}, gives
\[
 f_{\KK,1}=qA+O_R(q^3),\quad f_{\KK,2}=q^2B+O_R(q^4),
 \quad \dot f_{\KK,p,1}=qA_p^{\rm act}+O(q^3),\quad
 \dot f_{\KK,p,2}=q^2B_p^{\rm act}+O(q^4).
\]
Here scaled core velocities and reactions are uniformly bounded
by FWC's sourcewise equations and its moment bounds. Thus the
derivative expansion is justified directly, not by differentiating
a static \(O(q^4)\) symbol. It follows that
\begin{equation}
 q^{-2}D_p^{\KK,\rm self}
       =(A-\zeta)A_p^{\rm act}+B_p^{\rm act}+O(q^2).
 \label{an07rtav1:corealgebra}
\end{equation}

For explicit normalization, the physical leading source velocities
used for comparison are
\begin{align*}
 2F_{x1}&=-(1-M)x,&2F_{y1}&=-(1-M)y-Y-K_w,\\
 2F_{x2}&=\rho\phi(\rho x)-\rho\mathcal S(x)
                  +\lambda[V+(1-N)y]\Phi(\rho x),&
 2F_{y2}&=\rho(1-N)\Kfun(\rho x),
\end{align*}
with leading reactions \(h_1^0=1-M\), \(h_2^0=0\).
The moments and positive measures are FWC's physical ones;
\(Y=\int_\KK y\), \(V=\int_S v\),
\(K_w=\int_S\Kfun(u)\), and
\(\mathcal S(x)=\int_\KK\Ffun(\rho x,\rho\widetilde x)\).
Define \(A_p^0,B_p^0\) by replacing the velocities and reactions
in the preceding two integrals by \(F_{xp},F_{yp},h_p^0\).
Then \(d_p^{\rm lead}=(A-\zeta)A_p^0+B_p^0\) is precisely
the general leading-form rate used in CWI \texttt{rate}.
In particular the \(1-M\) reaction has not been dropped.
At \(M=1\) its first source reduces to
\(-A(Y+K_w)/2\), agreeing with P \texttt{probe-pop}.

FWC gives source velocity errors \(O(q^{1/10})\), and reaction
errors \(O(q^2+\mu_\BB)\) for source 1 and \(O(q^2)\) for source 2.
All coefficients in the two integrals are bounded on the core
chart, and \(\mu_\BB=O(q^{11/10})\). Subtracting them in
\eqref{an07rtav1:corealgebra} proves \eqref{an07rtav1:coreprecision}.
This uses physical \(M\), signed moments and sourcewise errors;
a total-field estimate alone would not supply this result.
\end{proof}

\begin{lemma}[Full background source charge]
\label{an07rtav1:backgroundrate}
Uniformly on \(\mathcal P_*\times J\),
\begin{equation}
 |D_p^\tau-D_p^{\KK,\rm self}|
                       \le C_{\rm bg}q^2\epsilon_{\rm bg}.
 \label{an07rtav1:backgroundprecision}
\end{equation}
\end{lemma}
\begin{proof}
Put \(T_0=\Sigma\cup\BB\). FWC gives
\(\mu_{T_0}\le Cq^{\gamma_*}\),
\(\mathcal Q_{T_0}\le Cq^2\epsilon_{\rm bg}\), and all
matrix-value/alignment hypotheses of CWI \texttt{ratecharge}
on the continued window. Hence
\begin{gather*}
 |f_{T_0}|\le\sqrt{\mu_{T_0}\mathcal Q_{T_0}},\qquad
 |\xi_q\cdot\dot f_{T_0,p}|\le
                  C(\mathcal Q_{T_0}+q^2L^2\mu_{T_0}),\\
 |\dot f_{T_0,p}|\le
                  C(\sqrt{\mu_{T_0}\mathcal Q_{T_0}}+qL\mu_{T_0}).
\end{gather*}
Also \(|f_\KK|+|\dot f_{\KK,p}|\le Cq\) by the reached
core chart and the preceding source calculation.
Primary inactivity leaves exactly
\[
 D_p^\tau-D_p^{\KK,\rm self}
 =f_{T_0}\cdot\dot f_{\KK,p}
   -(\xi_q-f_\KK)\cdot\dot f_{T_0,p}
   +f_{T_0}\cdot\dot f_{T_0,p}.
\]
Insert the displayed bounds; use
\(L^2\mu_{T_0}\le C\epsilon_{\rm bg}\) and Young's inequality,
in particular \(L\sqrt{\mu_{T_0}\epsilon_{\rm bg}}
\le C\epsilon_{\rm bg}\). This proves the result with every
output and rate cross term retained. No outside sign or chart
has been assumed.

At a possible background gate hit, changing the ReLU selector
within \([0,1]\) changes only its normal term in
\eqref{an07rtav1:exactrate}. There
\(|a_\psi\cdot\xi_q|\le|\theta-\psi|=O(qL)\) and the
source angular velocity is \(O(qL)\), since the input angle
is a bounded \(q\)-distance from an axis. The projected change
is at most \(Cq^2L^2\mu_{T_0}\), and the unprojected change
at most \(CqL\mu_{T_0}\). Thus the same bounds hold for the
standard selector, or any such selector. This proves a bound
on source formulas, not classical differentiability at an atom.
\end{proof}

\begin{theorem}[A-rate-transfer on the learned windows]
\label{an07rtav1:ratetransfer}
Let \(C_{\rm dist}\) be FWC's fixed-window constant in
\(\dist\le C_{\rm dist}q^\chi\), and set
\begin{equation}
 \delta_J(q)=C_{\rm core}(q^{1/10}+q^2)
      +C_{\rm bg}\epsilon_{\rm bg}
      +64R^2(1+4/g_*)C_{\rm dist}q^\chi.
 \label{an07rtav1:deltaJ}
\end{equation}
Then \(\delta_J(q)=O(q^{1/20})\to0\) and, for \(p=1,2\),
\begin{equation}
 \sup_{(\rho,\zeta)\in\mathcal P_*,\ s\in J}
 |q^{-2}D_p^\tau(\tau_h+s,\xi_q)-d_{p,\rho,\zeta}(s)|
                       \le\delta_J(q).
 \label{an07rtav1:ratelimit}
\end{equation}
This is uniform for the exact selector-defined source formulas;
it holds for classical source derivatives wherever defined and
in the almost-everywhere time sense. There is no core gate hit.
If
\begin{equation}
 \delta_J(q)<\min\{p_*/8,\gamma_0/16\},
 \label{an07rtav1:ratethreshold}
\end{equation}
then
\begin{align}
 D_1^\tau&\le-\tfrac12p_*q^2,&
 D_2^\tau&\ge\tfrac12p_*q^2&&\text{on }J,\notag\\
 D^\tau&\le-\tfrac58\gamma_0q^2&&&\text{on }J_-,\notag\\
 D^\tau&\ge\tfrac58\gamma_0q^2&&&\text{on }J_+.
 \label{an07rtav1:originalmargins}
\end{align}
\end{theorem}
\begin{proof}
Every hypothesis of CWI \texttt{rate} is matched: positive
physical measures; \(M,N\le2\); bounded strong support and
comparator; finite weak first moments; same \(\rho\); and
\(\zeta-X_*\ge g_*/4\) from WIT. It compares the leading-form
rate at the actual state with the coherent rate, with error
\(64R^2(1+4/g_*)\dist\). Apply it after
Lemmas \ref{an07rtav1:corerate} and
\ref{an07rtav1:backgroundrate}. Since \(\gamma_*>\chi\),
\(\epsilon_{\rm bg}=o(q^\chi)\), proving the stated order.
Each source error is below \(p_*/8\); their sum is below
\(\gamma_0/8\) under \eqref{an07rtav1:ratethreshold}.
Subtract from the exact limiting coefficients in
\eqref{an07rtav1:limitingmargins} to obtain
\(5/8-1/8=1/2\) for sources and
\(3/4-1/8=5/8\) for total rates.
Absolute continuity and the exact source chain rule identify
the sum with \(dE/d\tau\) almost everywhere; no derivative
of an error-transfer remainder is taken.
\end{proof}

\section{Learning of both concepts, including exact cluster losses}
\label{an07rtav1:learningsection}
WIT's \texttt{staticproposition} cannot be applied verbatim:
the entire original state is not supported in two fixed compact
scaled charts. The selected weak chart can grow like
\(q^{-603/640}\), and \(P_q,\Sigma,\BB\) remain present.
The following direct estimate replaces that hypothesis.
Set \(Q_1=1+q^2\), \(Q_2=\rho^2+q^2\), and
\[
 \mathsf m_p=Q_p^{-1}\E_{P_p}[X_pf_p(X)],\qquad
 \mathcal L_p=\tfrac12\E_{P_p}|f(X)-X|^2.
\]
In particular \(\mathsf m_2=c\) is the exact response,
not a name for the physical selected mass.

\begin{lemma}[Full-output learning estimate with weak moments]
\label{an07rtav1:learningestimate}
On \(J\), with \(M_1=M\), \(M_2=N\), there are fixed constants
\(C_{\rm out},C_{\rm resp},C_{\rm loss}\) such that
\begin{align}
 \|f-e_pM_pX_p\|_{L^2(P_p)}&\le C_{\rm out}q^{1/4},\notag\\
 |\mathsf m_p-M_p|&\le C_{\rm resp}q^{1/4},\notag\\
 \sqrt{2\mathcal L_p}&\le
       |1-M_p|\sqrt{Q_p}+C_{\rm loss}q^{1/4}.
 \label{an07rtav1:learningbounds}
\end{align}
All original populations are included.
\end{lemma}
\begin{proof}
FWC supplies bounded \(M,N\), a bounded strong chart, and
\(\int_S(1+|u|+|v|)\dd\mu\le C\).
It also gives
\(\mu_{P_q}+\mu_\Sigma+\mu_\BB=O(q^{1/4})\)
throughout the fixed window, by its fixed-label mass comparison.

For a core unit feature on \(P_1\), Lipschitz continuity of
the positive part gives
\[
 \|a_\psi(a_\theta\cdot X)_+-e_1X_1\|_2
 \le (|a_\theta-e_1|+|a_\psi-e_1|)\|X\|_2
                       +\|(-X_1)_+\|_2\le C_Rq.
\]
On \(P_2\) its feature norm is at most
\(q+q|x|\sqrt{Q_2}\le C_Rq\).
For a selected weak feature on \(P_1\), the norm is at most
\(q|u|\sqrt{Q_1}+q\). On \(P_2\), compare with \(e_2X_2\):
\[
 \|a_\psi(a_\theta\cdot X)_+-e_2X_2\|_2
 \le Cq(|u|+|v|)+\|(-X_2)_+\|_2
 \le Cq(1+|u|+|v|).
\]
The bounds \(|a_w-e_2|\le q|u|\),
\(|b_w-e_2|\le q|v|\) need no fixed weak supremum.
The Gaussian negative parts are bounded by \(q|Z|\), not
declared absent. Integrating these feature estimates against
the positive current measures costs \(O(q)\).
For each omitted cohort \(A\), the exact inequality
\(|f_A(X)|\le\mu_A|X|\) gives an unscaled \(L^2\) cost
\(C\mu_A\). This is legitimate at the present order-one
learning scale; it is not the \(q^2\)-precision probe argument.
Minkowski proves the first line of \eqref{an07rtav1:learningbounds}.
The response bound follows by Cauchy--Schwarz and
\(\sqrt{Q_p}\ge13/20\). Finally the unused target coordinate
has norm \(q\); absorb it into the last constant.
\end{proof}

\begin{proposition}[Responses and the fixed loss ratio]
\label{an07rtav1:learning}
For all sufficiently small \(q\), uniformly on \(K_\rho\times J\),
\begin{gather}
 M=1+O(q^{11/10}),\quad N=n_\rho+O(q^{1/8}),\quad
 c=N+O(q^{1/4}),\notag\\
 \tfrac34\le M\le\tfrac54,\qquad
 \tfrac14+\tfrac\sigma4\le N\le\tfrac54,\qquad
 \mathsf m_1\ge\tfrac12,\quad\mathsf m_2\ge\tfrac14,
 \label{an07rtav1:learningresponse}\\
 \mathcal L_p\le\frac{169}{196}\mathcal L_p(0),
                    \qquad p=1,2.
 \label{an07rtav1:lossratio}
\end{gather}
Thus both exact cluster losses have relative reduction at least
\(27/196\). No entire-state compact-chart hypothesis is retained.
\end{proposition}
\begin{proof}
The first line is FWC. WIT's mass margin and the uniform errors
give the two mass intervals if
\(C_Mq^{11/10}\le1/4\) and
\(C_Nq^{1/8}\le\min(1/4,\sigma/4)\).
For responses choose in addition
\(C_{\rm resp}q^{1/4}\le\min(1/4,\sigma/4)\).
Then \eqref{an07rtav1:learningbounds} gives the claimed thresholds.
Both masses are in \([1/4,5/4]\), so \(|1-M_p|\le3/4\).
If
\begin{equation}
 C_{\rm loss}q^{1/4}\le\frac{13}{320},
 \qquad S_0=q^{10}\le\tfrac12,
 \label{an07rtav1:learnthreshold}
\end{equation}
then \(C_{\rm loss}q^{1/4}/\sqrt{Q_p}\le1/16\) and
\(\sqrt{2\mathcal L_p}\le(13/16)\sqrt{Q_p}\).
The exact isotropic initial output is \(f_0(X)=S_0X/4\), hence
\[
 \sqrt{2\mathcal L_p(0)}
 =(1-S_0/4)\sqrt{Q_p+q^2}\ge(7/8)\sqrt{Q_p}.
\]
The squared ratio is \((13/14)^2=169/196\), with no limiting
loss substituted for the actual initialized one.
\end{proof}

\section{AN8: initialized Theorem A on the learned rectangle}
\label{an07rtav1:assembly}
All constants below are analytic, not numerical enclosures.
First keep the entire WPT/RB--WCS entry order and \(n\).
Freeze WIT's original rectangle, witnesses and their margins.
Choose \(T_J\), the FWC continuation constants and radius \(R\),
then \(L_F,C_{\rm sup},C_{\rm core},C_{\rm bg},C_{\rm dist}\)
and the direct learning constants. Let \(C_{\rm err}\) satisfy
FWC's normalized error-transfer bound on \(J\).
Choose \(q_0>0\) so that every \(0<q\le q_0\) satisfies all
source smallness thresholds and, explicitly,
\begin{gather}
 C_{\rm sup}q^\chi\le g_*/8,\qquad
 \delta_J(q)<\min(p_*/8,\gamma_0/16),\qquad
 C_{\rm err}q^\chi\le\gamma_0r/16,\notag\\
 C_Mq^{11/10}\le1/4,\quad
 C_Nq^{1/8}\le\min(1/4,\sigma/4),\quad
 C_{\rm resp}q^{1/4}\le\min(1/4,\sigma/4),\notag\\
 C_{\rm loss}q^{1/4}\le13/320,\qquad
 q^{10}\le1/2,\quad q\le1/2,\quad 2qR<\pi/2.
 \label{an07rtav1:q0}
\end{gather}
Such a \(q_0\) exists because all constants are fixed and all
the displayed errors vanish. It can also be decreased so that
\(\sup_J\dist\le\varepsilon_*\); this does not license direct
application of WIT's exact-leading-flow theorem to original dynamics.
No numerical value of \(q_0\) is asserted.

\begin{theorem}[Initialized learned signed competition and reversal]
\label{an07rtav1:theoremA}
For the original aligned-balanced isotropic initialization with
\(S_0=q^{10}\), every \(0<q\le q_0\), and every
\((\rho,\zeta)\in\mathcal P_*\), the common centered windows
\eqref{an07rtav1:windows} occur at normalized times \(\tau_h+J\)
and satisfy all of the following.
\begin{enumerate}
\item Both concepts are learned in the exact response/loss sense:
 \(\mathsf m_1\ge1/2\), \(\mathsf m_2\ge1/4\) and
 \(\mathcal L_p\le(169/196)\mathcal L_p(0)\), \(p=1,2\).
 The physical mass conclusions in
 \eqref{an07rtav1:learningresponse} also hold.
\item The original source rates satisfy
 \(D_1^\tau\le-p_*q^2/2<0<D_2^\tau\), with
 \(D_2^\tau\ge p_*q^2/2\), on \(J\).
 The total rate is at most \(-5\gamma_0q^2/8\) on \(J_-\)
 and at least \(5\gamma_0q^2/8\) on \(J_+\).
 These are the exact selector source observables; the corresponding
 classical derivatives and total derivative satisfy the same
 inequalities wherever defined and almost everywhere in time.
\item The continuous original probe error has an interior minimizer
 \(s_{\min,q}\in J^\circ\), possibly depending on the parameters,
 such that
 \begin{align}
 E(\tau_h+s_*-r)-E(\tau_h+s_{\min,q})
                &\ge c_{\rm rev}q^2,\notag\\
 E(\tau_h+s_*+r)-E(\tau_h+s_{\min,q})
                &\ge c_{\rm rev}q^2,\qquad
 c_{\rm rev}=\gamma_0r/4>0.
 \label{an07rtav1:rebound}
 \end{align}
 No uniqueness or transverse-zero claim is made.
\item For each \(\rho\in K_\rho\), these conclusions hold on the
 off-concept-cone angular sector
 \(\arg\xi_q=-\pi/2+qI_\zeta\), of exact width
 \(2\delta_\zeta q\).
\end{enumerate}
\end{theorem}
\begin{proof}
The response/loss conclusions are Proposition
\ref{an07rtav1:learning}; the source and total-rate margins are
Theorem \ref{an07rtav1:ratetransfer}. For the rebound use only
the continuous error transfer, so no rate differentiation at
background gate atoms is involved.
Each of \(J_-\), \(J_+\) has length \(r/2\). Integrating the
limiting margins gives a drop across \(J_-\) and a rise across
\(J_+\), each at least \(3\gamma_0r/8\) in scaled error.
FWC's two endpoint errors cost at most
\(2C_{\rm err}q^\chi\le\gamma_0r/8\). Thus the exact original
drop and rise are each at least \(\gamma_0rq^2/4\).
The minimum over the compact \(J\) exists. Comparing the left
endpoint with the end of \(J_-\), and the right endpoint with
the start of \(J_+\), excludes both endpoints as minimizers
and proves \eqref{an07rtav1:rebound}.
Finally \(I_\zeta\) has length \(2\delta_\zeta>0\); for small
\(q\) its probes have positive first coordinate and negative
second coordinate, outside the positive cone of the two concept
means. The angular map has derivative \(q\), proving the width.
\end{proof}

\paragraph{Physical time and rates.}
The physical windows are
\((\tau_h+J)/\mu_1^2\) and
\((\tau_h+J_\pm)/\mu_1^2\). Every normalized source or total-rate
bound is multiplied by \(\mu_1^2\):
\(D_p^{\rm phys}=\mu_1^2D_p^\tau\).
The error rebound and angular sector retain the factors \(q^2\)
and \(q\), respectively. The center uses the original analytic
clock, not a fitted time or the \(O(1)\)-accuracy value of \(t_b\).

\section{Initialized persistence on a diverging, unspecified horizon}
\label{an07rtav1:persistencesection}
\begin{corollary}[Initialized persistence after common clearing]
\label{an07rtav1:persistence}
There is a fixed finite \(S_{\rm clr}>s_*+r\), common to all
\(\mathcal P_*\), and there are functions \(T_q\to\infty\),
\(\varepsilon_q\to0\) as \(q\downarrow0\), such that for the
same original initialized family and center,
\begin{equation}
 \sup_{\substack{(\rho,\zeta)\in\mathcal P_*\\
                   S_{\rm clr}\le s\le S_{\rm clr}+T_q}}
 \left|E_{\xi_q}(\tau_h+s)-\tfrac12\right|
                       \le\varepsilon_q q^2.
 \label{an07rtav1:persistencebound}
\end{equation}
No rate of growth of \(T_q\) is asserted. This is not a
forever-in-time claim at a fixed positive \(q\).
\end{corollary}
\begin{proof}
The projection \(K_\rho\) is a compact interval in \((0,1)\).
WIT \texttt{branchsmooth} supplies continuous selected branches
on finite windows; the uniform forward growth bound in P, with
bounded entry on this compact set, supplies LP
\texttt{uniform-clear}'s other input. Take upper offset
\(\zeta_+=\max I_\zeta\). As already assembled in LPR
\texttt{combinedtheorem}, one common \(S_{\rm clr}>s_*+r\)
has \(X_\rho(s)>\zeta_+\) for all later \(s\).
P \texttt{atomicerror}, equivalently LPR's error formula, then
gives \(\mathcal E_{\rho,\zeta}(s)=0\) exactly there.
This is gate clearing, not merely a nonpositive benefit estimate.

For every finite \(T>S_{\rm clr}\), FWC
\texttt{errortransfer,\allowbreak early} gives the full original convergence
uniformly on \(K_\rho\times I_\zeta\times[S_{\rm clr},T]\).
Its family is the actual initialization \(S_0=q^{10}\),
and LP's required physical centering is
\(t_h^{\rm phys}=\tau_h/\mu_1^2\). Thus every hypothesis of
LP \texttt{diagonal-transfer} holds on exactly the same compact
parameter and offset sets. Apply that proposition to obtain
\eqref{an07rtav1:persistencebound}.
Its uniform threshold sequence can be chosen below \(q_0\);
for each integer \(j\), the compact-window error on
\([S_{\rm clr},S_{\rm clr}+j]\) is at most \(1/j\) once
\(q\le q_j\), with \(q_{j+1}\le q_j/2\).
The proposition then takes \(T_q=j(q)\), \(\varepsilon_q=1/j(q)\).
This records its uniformity and gives no new horizon-growth rate.
\end{proof}
In physical time the interval in
\eqref{an07rtav1:persistencebound} is
\([ (\tau_h+S_{\rm clr})/\mu_1^2,
    (\tau_h+S_{\rm clr}+T_q)/\mu_1^2]\).
In particular the original benefit \((1/2-E)_+\) is at most
\(\varepsilon_q q^2\) there, while the stronger absolute-error
statement is what the corollary proves.

\section{Dependency ledger, final scope and S6 handoff}
\label{an07rtav1:handoff}
The direct proof dependencies of this increment are:
\begin{itemize}\raggedright
\item FWC \texttt{main,\allowbreak primitive,\allowbreak charts,\allowbreak masslogistic,\allowbreak roughmass,\allowbreak 
corefield,\allowbreak distbound,\allowbreak background,\allowbreak errortransfer,\allowbreak early};
\item CWI \texttt{exactrate,\allowbreak ratecharge,\allowbreak rate,\allowbreak time}; PSE's
\texttt{matrix,\allowbreak conditional} enter through FWC's proved continuation;
\item P \texttt{leading,\allowbreak pop-sources,\allowbreak probe-pop,\allowbreak atomicerror,\allowbreak branch}
and the selected-orbit forward-growth result;
\item WIT \texttt{uniformtheorem,\allowbreak windows,\allowbreak rectangle,\allowbreak branchsmooth,\allowbreak 
generalrates,\allowbreak epsdef}, with \texttt{populationtransfer,\allowbreak 
staticproposition,\allowbreak clockalignment} audited in their exact scopes;
\item LP \texttt{uniform-clear,\allowbreak compact-transfer,\allowbreak diagonal-transfer}
and LPR \texttt{combinedtheorem}; H2 Sections 3.1--3.2 and AN8
specify the consumers.
\end{itemize}
No foundational, entry or passage argument is re-proved in this
assembly. The inherited initialized chain, whose conclusions are
encapsulated in FWC/PSE, is recorded for review:
\begin{itemize}
\item Exact model and target/leading identities: F and P.
\item Entry: PROFILE, OWPT, MOMC, Q2/Q2E, Q2EP, Q2EC and SEP.
\item Passage and budgets: LOC, CLOCK, TAILREP, BULK, RCET, LCRC,
CFR, DECG, GESG, S2CP, CFA, Q3SB, PISB, LSB, WPT and RB.
\item Matching and selection: WEAK, WBR, MP, SRR and WCS;
then CWI, PSE and FWC supply the compact-window interface.
\end{itemize}
Obstruction-only or superseded contracts are not substituted for
any of these completed conclusions. The companion review records
the exact repository paths and blob hashes read for this assembly.

The proved initialized parameter family is precisely
\(S_0=q^{10}\), \(0<q\le q_0\),
\((\rho,\zeta)\in\mathcal P_*\), with the analytic widths
inherited from WIT. It contains \((2/3,5\pi/9)\) in the
interior of its ratio/offset rectangle. It is not an initialized
interval of powers \(\kappa\), and no numeric lower bound for
\(q_0\) or for the local parameter widths is claimed.
In particular the canonical \((q,S_0)=(0.05,2\cdot10^{-4})\)
run is not covered; no claim is made for a rounded offset merely
because it is numerically close to \(5\pi/9\).

This completes the initialized learned-reversal and persistence
assembly in the stated deep regime. The original-flow B-core
lag/exposure and producer-attribution decomposition, including
accuracy relative to a small lag cancellation, remains S6.
The signed source bounds here do not supply that decomposition.
The non-small bridge and canonical regimes remain S7/S8.
There is no quantitative persistence horizon, no finite-width
claim and no assertion of unique reversal time.

For S6 the available handoff is the same matched clock; fixed-window
physical masses, current moments and pointwise sourcewise core
velocities/reactions; strict core probe gap on \(J\); the proved
source margins; and full-background source charges at the
\(q^2\epsilon_{\rm bg}\) scale. S6 must still identify and control
the exact lag/exposure terms in their required relative precision.
No checkpoint merge or change to any prior file is authorized or made.
\end{document}
